\documentclass[10pt]{beamer}
\usetheme{metropolis}

\usepackage{amsmath,amssymb}
\usepackage{mathtools}
\usepackage{xeCJK}

\setmainfont{Adobe 楷体 Std}


% 1. 开启标题背景色填充
\metroset{block=fill}


% 3. 增加标题与内容的间距
% theorem environments
\newtheorem{proposition}{Proposition}
\newtheorem{exercise}{Exercise}


% notation
\newcommand{\R}{\mathbb{R}}
\newcommand{\Q}{\mathbb{Q}}
\newcommand{\C}{\mathbb{C}}
\newcommand{\T}{\mathcal{T}}
\newcommand{\B}{\mathcal{B}}

\title{Elementary Topology}
\subtitle{Generalizing the concept of continuity to general sets of points}
\date{September 15, 2026}
\author{LXJ}

\begin{document}

\maketitle

%--------------------------------------------------------------------
\begin{frame}{Motivation: from calculus to topology}
	\begin{block}{Recall (continuity in $\R^N$)}
		A map $f:\R^N \to \R^M$ ($N=M=1$, for e.g.) is \emph{continuous} if
		\[
			\forall\, x_0 \in \R^N:\qquad \lim_{x \to x_0} f(x) = f(x_0),
		\]
		which is equivalent to the following description:
		\[
			\text{for every open set } U \subset \R^M,\quad f^{-1}(U) \text{ is open in } \R^N.
		\]
	\end{block}
	\begin{alertblock}{Key}
		To define continuity, we need the notion of \alert{open sets}.
	\end{alertblock}
\end{frame}

%====================================================================
\section{Topological Spaces}

\begin{frame}{Topological structure}
	\begin{definition}[Topological structure]
		Let $S$ be a non-empty set. A \emph{topological structure} $\T$ on $S$ is a
		collection of subsets of $S$ such that
		\begin{enumerate}
			\item $\varnothing,\ S \in \T$;
			\item if $U_\alpha \in \T$ for $\alpha \in I$ ($I$ being some index set),
			      then $\displaystyle\bigcup_{\alpha \in I} U_\alpha \in \T$;
			\item if $U_1, \dots, U_n \in \T$, then
			      $\displaystyle\bigcap_{i=1}^{n} U_i \in \T$.
		\end{enumerate}
		The pair $(S,\T)$ is called a \emph{topological space}.
	\end{definition}
	The elements of $\T$ are called \emph{open sets}.
\end{frame}

\begin{frame}{Closed sets}
	\begin{definition}[Closed set]
		The complement of an open set is called a \emph{closed set}.
	\end{definition}
\end{frame}

%====================================================================
\section{Continuous Maps}

\begin{frame}{Continuous maps between topological spaces}
	\begin{definition}[Continuous maps between topological spaces]
		A map $f: X \to Y$ between topological spaces $X$ and $Y$ is
		\emph{continuous} if for every open set $V \subset Y$, the preimage
		$f^{-1}(V)$ is an open set of $X$.
	\end{definition}
	\begin{proposition}[Composition of continuous maps]
		If $f: X \to Y$ and $g: Y \to Z$ are continuous maps, then
		$g \circ f: X \to Z$ is continuous.
	\end{proposition}
	\begin{proposition}[Continuity defined via closed sets]
		A map $f: X \to Y$ is continuous $\iff$ the preimage of every closed set
		is closed.
	\end{proposition}
\end{frame}

\begin{frame}{Open maps, closed maps, homeomorphisms}
	\begin{definition}[Open maps and closed maps]
		A map $f: X \to Y$ is called an \emph{open map} if the image
		$f(U) \subset Y$ is open for every open set $U \subset X$;\\
		$f: X \to Y$ is called \emph{closed} if $f(V) \subset Y$ is closed for
		every closed $V \subset X$.
	\end{definition}
	\begin{definition}[Homeomorphism (同胚)]
		A bijection (双射) $f: X \to Y$ is a \emph{homeomorphism} if both $f$ and
		$f^{-1}$ are continuous. Topological spaces $X$ and $Y$ are then said to
		be \emph{homeomorphic}.
	\end{definition}
\end{frame}

%====================================================================
\section{Subspace Topology}

\begin{frame}{Subspace topology}
	\begin{definition}[Subspace topology]
		Let $(S,\T)$ be a topological space and $A \subset S$ be a subset. Then
		\[
			\T_A \coloneqq \{\, A \cap U \mid U \in \T \,\}
		\]
		is the \emph{subspace topology} on $A$.
	\end{definition}
	\begin{exercise}
		Verify that $\T_A$ defines a topological structure on $A$.
	\end{exercise}
	\begin{exercise}
		Define the natural inclusion $\iota: A \hookrightarrow S$. Show that
		$(A,\T_A)$ is a subspace structure $\iff$ $\iota$ is continuous.
	\end{exercise}
\end{frame}

%====================================================================
\section{Bases}

\begin{frame}{Basis of a topology}
	\begin{definition}[Basis of a topology]
		A subcollection $\B \subset \T$ is a \emph{basis} for $\T$ if
		$\forall\, U \in \T$ and $p \in U$, $\exists\, B \in \B$ such that
		\[
			p \in B \subset U.
		\]
		We also say that $\B$ \emph{generates} the topology $\T$, or that $\B$ is
		a basis for the topological space $S$.
	\end{definition}
	\begin{example}
		$\B = \{\, B(p,r) \mid p \in \R^n,\ r > 0 \,\}$ is a basis for the
		standard topology on $\R^n$.
	\end{example}
\end{frame}

\begin{frame}{Characterizations of a basis}
	\begin{proposition}
		$\B$ is a basis for $\T$ iff every $U \in \T$ can be written as
		\[
			U = \bigcup_{\text{some } B \in \B} B.
		\]
	\end{proposition}
	\begin{proposition}[Criterion]
		A collection $\B$ of subsets of a set $S$ is a basis for a topology iff
		\begin{enumerate}
			\item $S = \displaystyle\bigcup_{B \in \B} B$;
			\item $\forall\, B_1, B_2 \in \B$ and $p \in B_1 \cap B_2$
			      (if it exists), $\exists\, B \in \B$ such that
			      $p \in B \subset B_1 \cap B_2$.
		\end{enumerate}
	\end{proposition}
	\begin{proposition}
		Let $\B = \{B_\alpha\}$ be a basis for the topological space $S$, and let
		$A$ be a subspace of $S$. Then $\{B_\alpha \cap A\}$ is a basis for $A$.
	\end{proposition}
\end{frame}

%====================================================================
\section{First and Second Countability}

\begin{frame}{Second countability}
	\begin{proposition}
		Let
		$\B_{\mathrm{rat}} \coloneqq \{\, B(p,r) \mid p \in \Q^n,\ r \in \Q^{+} \setminus \{0\} \,\}$.
		Then $\B_{\mathrm{rat}}$ is a basis for $\R^n$.
	\end{proposition}
	\begin{definition}[Second countable]
		A topological space is said to be \emph{second countable} if it has a
		\emph{countable basis}.
	\end{definition}
\end{frame}

\begin{frame}{First countability}
	\begin{definition}[Neighborhood basis]
		Let $S$ be a topological space and $p \in S$. A \emph{neighborhood basis}
		at $p$ is a collection $\B = \{B_\alpha\}$ of neighborhoods of $p$ such
		that for every neighborhood $U$ of $p$, $\exists\, B_\alpha \in \B$ with
		\[
			p \in B_\alpha \subset U.
		\]
	\end{definition}
	\begin{definition}[First countable]
		$S$ is said to be \emph{first countable} if it has a countable
		neighborhood basis at every point $p \in S$.
	\end{definition}
\end{frame}

%====================================================================
\section{Separation Axioms}

\begin{frame}{Hausdorff and normal spaces}
	\begin{definition}[Hausdorff topological spaces]
		A topological space $S$ is called a \emph{Hausdorff space} if
		$\forall$ points $x, y \in S$, $\exists$ two \emph{disjoint open sets}
		$U, V$ such that $x \in U$, $y \in V$.
	\end{definition}
	\begin{definition}[Normal spaces]
		A Hausdorff space is called \emph{normal} if, given any two disjoint
		closed sets $F, G$, $\exists$ disjoint open sets $U, V$ such that
		\[
			F \subset U \quad \& \quad G \subset V.
		\]
	\end{definition}
	\begin{proposition}
		Any one-point set in a Hausdorff space is closed.
	\end{proposition}
	\begin{exercise}
		Any subspace of a Hausdorff space is also Hausdorff.
	\end{exercise}
\end{frame}

\begin{frame}{Examples}
	\begin{example}
		$\R^n$ is a Hausdorff topological space.
	\end{example}
	\begin{example}[The Zariski topology]
		Let $S = \C^n$ be the $n$-dimensional vector space over $\C$, endowed
		with the \emph{Zariski topology}. Every open set $U$ is of the form
		$S \setminus Z(I)$, where $I$ is a family of polynomials in $n$ variables
		and $Z(I)$ is the solution set of this family.
		\par\smallskip
		Any two such open sets intersect:
		\[
			U \cap V = S \setminus \bigl( Z(I) \cup Z(J) \bigr)
		\]
		is nonempty, since $Z(I) \cup Z(J)$ --- being the solution set of $I$ or
		$J$ --- cannot be the whole space. Hence the Zariski topology is
		\emph{not} Hausdorff.
	\end{example}
\end{frame}

%====================================================================
\section{Product Topology}

\begin{frame}{Product topology}
	The Cartesian product $A \times B = \{(a,b) \mid a \in A,\ b \in B\}$.
	\begin{definition}[Product topology]
		Given two topological spaces $X$ and $Y$, construct the collection $\B$
		consisting of all sets of the form $U \times V$ ($U, V$ open sets in $X$
		and $Y$). $\B$ generates a topology on $X \times Y$, called the
		\emph{product topology}.
	\end{definition}
	\begin{proposition}
		Let $\{U_i\}$ and $\{V_j\}$ be bases for $X$ and $Y$ respectively. Then
		$\{U_i \times V_j\}$ is a basis for $X \times Y$.
	\end{proposition}
	\begin{proposition}
		The product of two second countable topological spaces is second
		countable.
	\end{proposition}
	\begin{exercise}
		The product of two Hausdorff topological spaces is Hausdorff.
	\end{exercise}
\end{frame}

%====================================================================
\section{Compactness}

\begin{frame}{Covers and compact spaces}
	\begin{definition}[Cover, subcover]
		Let $S$ be a topological space. A collection of open sets
		$\{U_\alpha\}_{\alpha \in A}$ of $S$ is called a \emph{cover} (or
		\emph{open cover}) of $S$ if $S = \bigcup_{\alpha} U_\alpha$. If a
		subcollection $\{U_\alpha\}_{\alpha \in B \subset A}$ of
		$\{U_\alpha\}_{\alpha \in A}$ still covers $S$, the subcollection is
		called a \emph{subcover} of $S$.
	\end{definition}
	\begin{definition}[Compact topological space]
		A topological space $S$ is called \emph{compact} if every open cover of
		$S$ has a \emph{finite subcover}.
	\end{definition}
\end{frame}

\begin{frame}{Compactness: basic properties}
	\begin{exercise}
		A closed subset of a compact topological space is compact.
	\end{exercise}
	\begin{exercise}
		Every compact subset $K$ of a Hausdorff space $S$ is closed.
	\end{exercise}
	\begin{proposition}
		The image of a compact set under a continuous map is compact.
	\end{proposition}
\end{frame}

\begin{frame}{Maps from compact to Hausdorff spaces}
	\begin{proposition}
		A continuous map from a compact topological space $X$ to a Hausdorff
		topological space $Y$ is a closed map.
	\end{proposition}
	\begin{corollary}
		A continuous bijection $f$ from a compact topological space $X$ to a
		Hausdorff space $Y$ is a homeomorphism.
	\end{corollary}
	\begin{exercise}
		Prove the corollary above.
	\end{exercise}
\end{frame}

%====================================================================
\section{Connectedness}

\begin{frame}{Connected spaces and connected components}
	\begin{definition}[Connected topological spaces]
		A topological space $S$ is called \emph{connected} if there do
		\emph{not} exist two nonempty disjoint open subsets $U$ and $V$ such
		that $S = U \cup V$. A subset $A \subset S$ is called \emph{connected}
		if it is connected as a topological subspace.
	\end{definition}
	\begin{proposition}
		The image of a connected topological space $X$ under a continuous map
		$f: X \to Y$ is connected.
	\end{proposition}
	\begin{definition}[Connected components]
		Let $x$ be a point of $S$. The set $C_x$ --- the union of all connected
		subsets of $S$ containing $x$ --- is called the \emph{connected
			component} of $S$ containing $x$.
	\end{definition}
\end{frame}

%====================================================================
\section{Closure}

\begin{frame}{Closure}
	\begin{definition}[Closure]
		The \emph{closure} of $A$ in $S$, denoted by $\overline{A}$, is defined
		as the intersection of all closed sets containing $A$ (the smallest
		closed set containing $A$).
	\end{definition}
	\begin{proposition}[Local characterization of the closure]
		Let $A \subset S$. A point $p \in S$ lies in $\overline{A}$ iff every
		neighborhood of $p$ contains at least one point of $A$.
	\end{proposition}
	\begin{proposition}
		A subset $A \subset S$ is closed iff $A = \overline{A}$.
	\end{proposition}
	\begin{proposition}
		If $A \subset B$, then $\overline{A} \subset \overline{B}$.
	\end{proposition}
\end{frame}

%====================================================================
\section{Further Exercises}

\begin{frame}{Exercises and further reading}
	\begin{exercise}
		What is a \emph{paracompact} topological space?
	\end{exercise}
	\begin{exercise}
		Read the part on the \emph{quotient topology}.
	\end{exercise}
\end{frame}

\end{document}
